% HAND-WRITTEN fragment -- placed by the `order` list in
% configs/tb_anatop_tia_rprog_corners.json.
%
% Specification grading of the 2026-08-21 spec-evaluated runs. Numbers come
% from the run databases (plain SQLite; tables spec/specValue/resultValue):
%   simruns/tb_anatop_tia_rprog_specrun/castalia/tb_anatop_tia_rprog/maestro/
%     results/maestro/Ocean.0.rdb           (rp_code + rp_lin, all 5 corners)
%   simruns/tb_anatop_tia_rprog_mcspec/.../Ocean.0.rdb  (rp_mc, 200 samples)
% and from the summary_*.txt files beside them. The MC rows repeat
% tab_tiarprog_mc, which is regenerated from the same run.

The bench carries 14 pass/fail specifications, graded by the simulator against
every corner and every Monte Carlo sample
(Table~\ref{tab:tiarprog-spec}). The limits are screening bounds chosen from
the measured data: the endpoint windows allow the \(\pm 35\,\%\) poly
sheet-resistance span plus margin, so a violation flags a broken device or
model, not process spread. Twelve of the fourteen pass at every corner and
every sample.

The two exceptions are honest and bounded. \texttt{v\_top} exceeds its
\SI{2.45}{\volt} limit at the two slow-resistor corners
(\SI{2.553}{\volt} at \texttt{ss-res}, \SI{2.518}{\volt} at \texttt{ss 125C}):
at \(+35\,\%\) resistance the fixed \SI{500}{\nano\ampere} test current rails
the \SI{2.5}{\volt} supply at the top codes. That is the measurement drive
running out of compliance at the slow corner, not a device defect --- the same
compliance ceiling already documented as the \(\sim\SI{0.65}{\micro\ampere}\)
top-code operating limit --- and the \SI{250}{\nano\ampere} control run bounds
its effect on the reported resistances at \(\le 0.036\,\%\)
(Table~\ref{tab:tiarprog-corners}). \texttt{lsb\_err} fails its 0.5-LSB budget
at four of five corners (1.24 typical, 3.60 worst at \texttt{ss 125C}, passing
only at \texttt{ff -40C} with 0.41): this is the bottom-code switch-Ron bowl of
Figure~\ref{fig:tiarprog-lin-n0}, largest where the gates are slowest and
hottest. It is accepted as documented --- the chord-resistance error
\texttt{r\_lin\_err\_pct} stays below \(1\,\%\) everywhere
(0.12--0.53\,\%), and the fix, if bottom-code linearity ever matters, is wider
bypass gates. One bench note: the linearity sweep's drive source includes the
\texttt{itest} term (\(i = i_{\mathrm{test}} + f\cdot 0.8/R_{\mathrm{ideal}}\),
mirroring the netlist source exactly), so the chord is taken about the same
operating point the code sweep uses; against that corrected drive the typical
worst deviation reads 1.24 LSB where Table~\ref{tab:tiarprog-lin} reported
1.26.

% Hand-written table; values transcribed from the rdb spec evaluation above.
\begin{table}[htbp]
  \centering
  \caption[{Specification grading of the programmable feedback resistor: 14
  simulator-graded limits.}]{Specification grading: the 14 limits carried by
  the bench and the verdict. Corner rows list the limit and the measured
  extremes across all five corners of Table~\ref{tab:TiaRprog-corners}; the
  endpoint and step windows are screening bounds (\(\pm 35\,\%\) corner span
  plus margin). Monte Carlo rows list the nominal (mean) and standard
  deviation over the 200 mismatch samples (typical process,
  \SI{25}{\celsius}), each graded against a \(\pm 5\,\%\) absolute-accuracy
  window about its typical value and the seam against zero.}
  \label{tab:tiarprog-spec}
  \small
  \begin{tabular}{llrrl}
    \toprule
    Specification & Unit & Limit & Measured & Verdict \\
    \midrule
    \multicolumn{5}{l}{\emph{Code sweep, 320 codes, five corners}} \\
    $R$ at $N{=}0$ & \si{\kilo\ohm} & 12--28 & 13.26--26.41 & pass \\
    $R$ at $N{=}319$ & \si{\mega\ohm} & 1.25--2.65 & 1.280--2.606 & pass \\
    Step, min & \si{\kilo\ohm} & $>1$ & 3.57--7.67 & pass \\
    Step, mean & \si{\kilo\ohm} & 3.8--8.2 & 3.97--8.09 & pass \\
    Non-monotonic steps & count & 0 & 0 & pass \\
    Error vs $18\,\mathrm{k}+6\,\mathrm{k}\cdot N$, max & \si{\percent} & $<55$ & 9.2--46.7 & pass \\
    $V(\mathtt{A})$ at top code & \si{\volt} & $<2.45$ & 1.89--2.55 & marginal$^{a}$ \\
    \addlinespace
    \multicolumn{5}{l}{\emph{Drive linearity at $N{=}0$, five corners}} \\
    Chord deviation, max & ADC LSB & $<0.5$ & 0.41--3.60 & fail$^{b}$ \\
    Secant-$R$ error & \si{\percent} & $<1$ & 0.12--0.53 & pass \\
    \addlinespace
    \multicolumn{2}{l}{\emph{Monte Carlo, 200 mismatch samples, typical}}
      & Nominal & Std.\ dev. & \\
    $R(N{=}0)$ & \si{\kilo\ohm} & 19.650 & 0.071 & pass, 200/200 \\
    $R(N{=}63)$ & \si{\kilo\ohm} & 396.452 & 0.361 & pass, 200/200 \\
    $R(N{=}64)$ & \si{\kilo\ohm} & 403.876 & 0.364 & pass, 200/200 \\
    $R(N{=}319)$ & \si{\mega\ohm} & 1.93212 & 0.00071 & pass, 200/200 \\
    Seam step $63{\to}64$ & \si{\kilo\ohm} & 7.424 & 0.504 & pass, 200/200 \\
    \bottomrule
    \multicolumn{5}{l}{\footnotesize $^{a}$\,Exceeded at \texttt{ss-res}
      (\SI{2.553}{\volt}) and \texttt{ss 125C} (\SI{2.518}{\volt}): test-drive
      compliance, not a device defect (see text).} \\
    \multicolumn{5}{l}{\footnotesize $^{b}$\,Fails at four of five corners;
      passes at \texttt{ff -40C} (0.41). Bottom-code switch-Ron bowl, accepted
      (see text).} \\
  \end{tabular}
\end{table}

% =====================================================================
% 2026-08-25 ADDITION -- in-loop noise.  HAND-WRITTEN; appended here rather
% than given its own fragment because TiaRprog.tex is maestro2tex output and
% its input list cannot be hand-extended.
%
% Source: ~/chips/castalia/simruns/bgladder_20260825/README.md Task B2, with
% tables/rprog_noise.tsv, rprog_noise/summary_rp_noise.txt and the A/B decks
% under rprog_noise/ab_idealRf/.  Bench castalia/tb_anatop_tia_rprog_noise/
% maestro test rp_noise; design pointer on castalia/tb_anatop_tia_amp_ab/
% TiaFbRprogTB, i.e. the real ladder inside the real TIA loop.  stb + ac +
% noise at 37 C, C_dl = 100 nF, R_u = 100 ohm, R_ct = 1 Mohm, C_f = 22 pF,
% section castalia_typ_25, noise 0.1 Hz - 1 MHz at 20 pts/dec.  Six corners =
% six gain codes.
%
% i_n is DERIVED, not asked of spectre: noise and ac share the frequency axis
% and I_WE carries mag=1, so i_n(f) = v_n(f)/|Z_T(f)| is exact.  The analysis's
% inType/inIsrc fields netlist to nothing in this state schema and getData("in")
% then errors, silently dropping the output.
%
%   N     R_f      i_n 0.1-10 Hz  i_n 0.1-1 kHz  ladder Johnson  ladder share
%   0     19.63k   1060.5 pA      3303.5 pA      934 fA/rtHz     0.0008 %
%   63    394.8k     81.4 pA      3005.2 pA      208             0.0065 %
%   127   777.1k     60.0 pA      3004.0 pA      148             0.0061 %
%   191   1.159M     53.6 pA      3003.6 pA      122             0.0051 %
%   255   1.542M     50.6 pA      3003.5 pA      105             0.0043 %
%   319   1.924M     48.9 pA      3003.4 pA       94             0.0037 %
%   i_n at 1 kHz is 132-133 fA/rtHz at EVERY code; 0.1-10 kHz is 26.7 nA at
%   every code.  PM 74.94 deg (N=63) - 92.86 deg (N=319), GM 21.65-21.79 dB,
%   Zt peaking <= 1.098 dB -- identical to the V-1x grid, so the noise run did
%   not perturb the operating point.
%   A/B against a single ideal resistor of the MEASURED value, standalone
%   spectre decks: N=0 1060.509 vs 1060.508 pA (-0.0001 %), N=63 81.372 vs
%   81.374 (+0.0020 %), N=319 48.845 vs 48.854 (+0.0179 %).
%   "Amplifier" here means everything that is not the feedback resistor: the
%   class-AB amplifier, the bias generator through it, and the electrode's own
%   R_u/R_ct.  noiseSummary() could not be driven headlessly, so the split is
%   not finer than that.
% One electrode and one process corner only (see note_tiarprog_scope).
% =====================================================================

\subsubsection{In-loop noise}\label{sss:tiarprog-noise}

Inside the transimpedance loop the ladder is not a noise contributor. Its own
Johnson noise is \SIrange{0.0008}{0.0065}{\percent} of the input-referred noise
power in the dc-amperometric band, four to five orders below the amplifier
(Table~\ref{tab:tiarprog-noise}), and replacing the whole switched ladder by a
single ideal resistor of the measured value changes the answer by
\SI{0.02}{\percent} at the worst code. The roughly three hundred open bypass
switches contribute nothing measurable. Loop stability is unchanged by the noise
run --- phase margin \SIrange{74.9}{92.9}{\degree}, gain margin
\SIrange{21.65}{21.79}{\decibel} --- so these are the operating points the
stability grid already reports.

\textbf{No noise figure for this block is meaningful without its band.} Over
\SIrange{0.1}{10}{\hertz} the input-referred current noise is
\SIrange{48.9}{1060}{\pico\ampere} depending on gain code. Over
\SI{0.1}{\hertz}--\SI{1}{\kilo\hertz} it is \SI{3.00}{\nano\ampere} at
\emph{every} code, and over \SI{0.1}{\hertz}--\SI{10}{\kilo\hertz}
\SI{26.7}{\nano\ampere} at every code. That is not a ladder defect: above the frequency at which the feedback
impedance exceeds the electrode impedance the noise gain is set by the electrode
capacitance, so the input-referred noise becomes independent of \(R_f\) --- which
is what the identical \(\SI{132}{\femto\ampere}/\sqrt{\si{\hertz}}\) spot
noise at \SI{1}{\kilo\hertz} across six decades of \(R_f\) shows. The same design is
\SI{49}{\pico\ampere} or \SI{27}{\nano\ampere} according to where the integral
stops, so the band edges are fixed design variables rather than tracked to each
code's own bandwidth --- letting them follow would make the slow, high-gain
codes report less noise merely for being slower.

Against the \SI{2}{\nano\ampere} input-referred budget in the amperometric band,
codes \(N \ge 63\) sit \(25\times\) to \(40\times\) under it and the bottom code
\(1.9\times\) under. At the top code the measured \SI{48.9}{\pico\ampere} is
\(17\times\) below the \SI{0.824}{\nano\ampere} current quantum, so the
resolution floor at maximum gain is quantisation, not noise.

\begin{table}[htbp]
  \centering
  \caption[{Input-referred current noise of the assembled transimpedance loop against gain code.}]{Input-referred current noise of the assembled transimpedance loop against gain code, \SI{37}{\celsius}, nominal process, at a \SI{100}{\nano\farad} double-layer capacitance and the \SI{22}{\pico\farad} compensation capacitor. Noise is referred through the measured transimpedance rather than by an input-referral option, which does not reach the netlist from this setup. The last column is the feedback resistor's own Johnson noise as a fraction of the noise power in the \SIrange{0.1}{10}{\hertz} band. The two bands differ because the noise gain above the feedback/electrode impedance crossover is set by the electrode, not by \(R_f\).}
  \label{tab:tiarprog-noise}
  \begin{tabular}{rrrrr}
    \toprule
    $N$ & $R_f$ & \multicolumn{2}{c}{Input-referred noise, rms} & Ladder share \\
    \cmidrule(lr){3-4}
     & & \SIrange{0.1}{10}{\hertz} & \SI{0.1}{\hertz}--\SI{1}{\kilo\hertz} & of noise power \\
    \midrule
    0   & \SI{19.63}{\kilo\ohm} & \SI{1060.5}{\pico\ampere} & \SI{3.304}{\nano\ampere} & \SI{0.0008}{\percent} \\
    63  & \SI{394.8}{\kilo\ohm} & \SI{81.4}{\pico\ampere}   & \SI{3.005}{\nano\ampere} & \SI{0.0065}{\percent} \\
    127 & \SI{777.1}{\kilo\ohm} & \SI{60.0}{\pico\ampere}   & \SI{3.004}{\nano\ampere} & \SI{0.0061}{\percent} \\
    191 & \SI{1.159}{\mega\ohm} & \SI{53.6}{\pico\ampere}   & \SI{3.004}{\nano\ampere} & \SI{0.0051}{\percent} \\
    255 & \SI{1.542}{\mega\ohm} & \SI{50.6}{\pico\ampere}   & \SI{3.003}{\nano\ampere} & \SI{0.0043}{\percent} \\
    319 & \SI{1.924}{\mega\ohm} & \SI{48.9}{\pico\ampere}   & \SI{3.003}{\nano\ampere} & \SI{0.0037}{\percent} \\
    \bottomrule
  \end{tabular}
\end{table}
